\section{Agentic \& Tool Use}

\subsection{当前 Agent 的核心瓶颈}

\textbf{谢天宝}：过去一年最大突破是从短时任务到多步长时任务。但\textbf{决策质量仍然很低}——观察 trajectory 会发现关键 decision 经常出错（Deep Research 的 report 不反映事实，Coding Agent 引入 bug 后反复 debug）。第二大问题是 \textbf{Context/Memory 管理}：当 multimodal agent 的 token 足够多时，很难在无损信息的前提下存储、管理和利用 memory，这导致 agent 有能力上限。Memory 的表示、管理和利用——本质上还是 2020--2023 年的研究方式（如 symbolic memory management），并未根本性改变。

\textbf{刘子伟}：三个关键点：
\begin{enumerate}
  \item \textbf{Long Sequence}——Tool Use 越实用，状态越长；end-to-end 解需要突破 long sequence 的 action space
  \item \textbf{Self Reflection}——每步出错概率随步数累积，强纠错/恢复能力是 real-world 落地的关键
  \item \textbf{技能迁移}——Computer Use 中学到的 skills 有没有可能 transfer 到 Physical World（大部分 robots 不会用工具）。人之所以快速上手电脑图标，是因为图标抽象了现实工具的元素——反过来也可能成立。
\end{enumerate}

\begin{importantbox}{曹玉：2025 年最 Impressive 的 Agentic 产品}
\textbf{软件}：Claude Code——强长程 agentic 能力，在 digital world 中做到了理想中 agentic AI 应做的事。由此可见 Anthropic 对 Agentic \& Tool Use 的理解达到了很高水位。整个产业因此发生变化，大部分 coding 服务商都把 Claude Code 作为核心服务之一。

\textbf{硬件}：豆包手机——可以自动刷抖音极速版赚钱并用微信转账，非常 impressive。但几天后很多功能被限制（不能打开微信、不能跨平台比价），暴露了\textbf{agentic AI 技术水平与国内厂商开放政策之间的矛盾}。

核心观察：在 Digital World 中 scaling 到 trillions of environments 很难。一种可能的 paradigm 是\textbf{直接把 agent 推向真实世界}（如 Claude Code 做软件、豆包手机做硬件），让世界给出 reward signal，而不是在模拟环境中无休止地训练。
\end{importantbox}

\textbf{薛福昭}：对 agentic 不是专家，但相信从 Tool Use $\to$ Web Agent $\to$ \textbf{Gaming}（Minecraft $\to$ 魔兽 $\to$ GTA 级 3A 游戏）$\to$ Physical AI 的路径。Gaming 是重要的中间步骤——3A 游戏提供了非常真实的 simulator 环境。

\subsection{本章小结}
Agent 已能处理长时任务但决策质量低、Memory 管理是硬瓶颈。Claude Code 是 2025 年 agentic 的标杆产品。真实世界的 reward signal 可能比模拟环境更有效。Gaming 是通向 Physical AI 的关键中间步骤。

%% ===== Continual Learning =====
